%File: main.tex
%====================================================================
% AAAI-style report for HC-RAG project.
% Single source file (AAAI requirement). Edit between the
%   "% === <SectionName> ==="
% markers; macros for numeric results live in the EDITABLE NUMBERS
% block so updates as new data arrives are one-line changes.
%====================================================================

\documentclass[letterpaper]{article}
\usepackage{aaai}
\usepackage{times}
\usepackage{helvet}
\usepackage{courier}
\usepackage{graphicx}
\usepackage{url}
\frenchspacing
\setlength{\pdfpagewidth}{8.5in}
\setlength{\pdfpageheight}{11in}

% --- PDF metadata (PLAIN ASCII ONLY; no LaTeX commands, no accents) ---
\pdfinfo{
/Title (Higher Criticism for Adaptive Retrieval in RAG Systems)
/Author (Lihu Zur and Noam Delbari)
/Keywords (Retrieval-Augmented Generation, Higher Criticism, Adaptive Retrieval, Information Retrieval, Large Language Models)
}

% --- Section numbering: 0=none, 1=sections, 2=subsections ---
\setcounter{secnumdepth}{2}

%====================================================================
% === Editable Numbers ===
% Update these macros when new experimental results arrive.
% Reference them in the body as e.g. \HCemHotpot
%====================================================================
% CrossEntityQA vs Top-K (k=15)
\newcommand{\CEqaEMHC}{0.751}
\newcommand{\CEqaEMTopK}{0.500}
\newcommand{\CEqaPrecHC}{0.711}
\newcommand{\CEqaPrecTopK}{0.670}
\newcommand{\CEqaRecHC}{0.572}
\newcommand{\CEqaRecTopK}{0.610}

% Musique vs Top-K (k=15)
\newcommand{\MusEMHC}{0.453}
\newcommand{\MusEMTopK}{0.561}
\newcommand{\MusPrecHC}{0.611}
\newcommand{\MusPrecTopK}{0.542}
\newcommand{\MusRecHC}{0.405}
\newcommand{\MusRecTopK}{0.550}
\newcommand{\MusFOneHC}{0.487}
\newcommand{\MusFOneTopK}{0.546}

% HotpotQA vs Adaptive-K and Adaptive-RAG (IRCoT)
\newcommand{\HotEMHC}{0.416}
\newcommand{\HotEMIRCoT}{0.446}
\newcommand{\HotEMAdaK}{0.338}
\newcommand{\HotPrecHC}{0.562}
\newcommand{\HotPrecIRCoT}{0.624}
\newcommand{\HotPrecAdaK}{0.281}
\newcommand{\HotRecHC}{0.559}
\newcommand{\HotRecIRCoT}{0.637}
\newcommand{\HotRecAdaK}{0.794}
\newcommand{\HotFOneHC}{0.541}
\newcommand{\HotFOneIRCoT}{0.604}
\newcommand{\HotFOneAdaK}{0.338}

% HotpotQA efficiency
\newcommand{\HotTokHC}{1{,}320.2}
\newcommand{\HotTokAdaK}{25{,}249.9}
\newcommand{\HotDocsHC}{10.9}
\newcommand{\HotDocsAdaK}{227.6}
\newcommand{\HotLatHC}{15.1}
\newcommand{\HotLatAdaK}{154.9}

% Holobench vs Adaptive-K
\newcommand{\HolEMHC}{0.374}
\newcommand{\HolEMAdaK}{0.326}
\newcommand{\HolPrecHC}{0.890}
\newcommand{\HolPrecAdaK}{0.531}
\newcommand{\HolRecHC}{0.606}
\newcommand{\HolRecAdaK}{0.685}
\newcommand{\HolFOneHC}{0.706}
\newcommand{\HolFOneAdaK}{0.497}
\newcommand{\HolTokHC}{5{,}849}
\newcommand{\HolTokAdaK}{20{,}234}
\newcommand{\HolDocsHC}{179.1}
\newcommand{\HolDocsAdaK}{638.4}
\newcommand{\HolLatHC}{1.46}
\newcommand{\HolLatAdaK}{11.49}

% Category-Key (44 country queries, 726 docs, BGE)
\newcommand{\CatHCF}{0.893}
\newcommand{\CatHCR}{0.852}
\newcommand{\CatHCP}{0.938}
\newcommand{\CatHCK}{9.7}
\newcommand{\CatCorrK}{0.835}

% Amazon Multi-Department (80K, 135 queries, OpenAI emb)
\newcommand{\AmMdHCF}{0.204}
\newcommand{\AmMdHCK}{34.1}
\newcommand{\AmMdTwentyF}{0.185}
\newcommand{\AmMdAvgF}{0.192}
\newcommand{\AmMdFiftyF}{0.187}
\newcommand{\AmMdHCTokF}{0.299}
\newcommand{\AmMdTwentyTokF}{0.261}

% Amazon Electronics / Book Categories
\newcommand{\AmElHCF}{0.192}
\newcommand{\AmElTwentyF}{0.185}
\newcommand{\AmBkHCF}{0.133}
\newcommand{\AmBkFiftyF}{0.128}

% Algorithm comparison on Amazon Multi-Dept (pool=1000)
\newcommand{\AlgHCF}{0.209}
\newcommand{\AlgBJF}{0.202}
\newcommand{\AlgBonfF}{0.022}
\newcommand{\AlgBHF}{0.074}
\newcommand{\AlgAdaKF}{0.138}
\newcommand{\AlgCARF}{0.110}

%====================================================================
% === Title and Author ===
%====================================================================
\title{Higher Criticism for Adaptive Retrieval\\in Retrieval-Augmented Generation}
\author{Lihu Zur \and Noam Delbari\\
Final Project Report\\
\texttt{lihu.zur@example.com} \quad \texttt{noam.delbari@example.com}\\
}

\begin{document}
\maketitle

%====================================================================
% === Abstract ===
%====================================================================
\begin{abstract}
\begin{quote}
Retrieval-Augmented Generation (RAG) pipelines typically retrieve a
fixed number of passages per query, wasting context budget on easy
questions and starving hard ones. We study \emph{Higher Criticism}
(HC), a statistical test for rare and weak signals, as a
query-adaptive cutoff mechanism layered on top of any dense retriever.
We integrate HC into three widely used RAG pipelines (Adaptive-K,
Adaptive-RAG, FlashRAG) and into our own end-to-end evaluation
pipeline, calibrate a per-corpus null distribution offline, and
compare HC against fixed Top-K, Adaptive-K, and the multi-step
Adaptive-RAG router across five evaluation datasets, with an
additional focused failure-analysis study on BEIR FiQA. We also
introduce \emph{CrossEntityQA}, a sparse multi-entity retrieval
benchmark in which only $\sim$2\% of the corpus is relevant per query.
On the sparse-relevance regime HC lifts Exact Match over fixed
Top-K (CrossEntityQA: $+50\%$ relative) and uses up to $95\%$ fewer
input tokens than Adaptive-K; it dominates Adaptive-K on long-context
database-style reasoning (HoloBench), but underperforms on
dependency-heavy multi-hop questions (MuSiQue vs.\ Top-K, HotpotQA
vs.\ multi-step IRCoT). On controlled and compound-query benchmarks
(Category-Key, Amazon Multi-Department / Electronics / Books) HC
also beats every fixed top-$k$ on F1 and dominates five alternative
adaptive cut-off rules (Berk--Jones, Bonferroni, BH, gap-based
Adaptive-K, clustering-based CAR). We characterise when HC helps,
when it fails, and outline a hybrid rerank-plus-HC direction.
\end{quote}
\end{abstract}

%====================================================================
% === Background ===
%====================================================================
\section{Background and Motivation}

Retrieval-Augmented Generation (RAG), introduced by
\citeauthor{lewis2020rag}~(\citeyear{lewis2020rag}), augments a Large
Language Model (LLM) with passages pulled from an external corpus,
typically via dense similarity search over an
approximate-nearest-neighbour index. The
dominant interface is \emph{Top-K}: fix an integer $k$ and pass the
$k$ most similar passages to the generator. This choice is convenient
but blunt --- a single $k$ cannot simultaneously serve a fact-lookup
query that needs one passage, a multi-hop question that needs three,
and a database-style aggregation that needs dozens. Too small a $k$
discards evidence; too large a $k$ inflates context length, latency,
and dollar cost, and degrades generation quality through distractor
noise.

\subsection{Existing Adaptive Approaches}
Two families of remedies have been proposed. \emph{Adaptive-K}
\cite{adaptivek2025} chooses $k$ per query by locating the largest gap
in the sorted similarity scores. \emph{Adaptive-RAG}
\cite{jeong2024adaptiverag} trains a small classifier that routes each
query to one of \{no retrieval, single-step retrieval, iterative
IRCoT\}. The first is cheap but brittle when score distributions are
flat; the second is accurate on multi-hop questions but pays for it
with many retriever and generator calls.

\subsection{Retrieval as Sparse Signal Detection}
\label{sec:sparse}
A query-aware retriever projects every document in the corpus onto a
similarity score against the query. The vast majority of those scores
come from \emph{irrelevant} passages and together form a background
``null'' distribution. Only a small, query-dependent number of
relevant passages (the \emph{signals}) are pushed above this
background. Choosing $k$ is therefore not really a ranking problem
but a \emph{detection} problem: how many of the top scores are too
unusual to have been produced by the irrelevant majority alone? This
is exactly the rare/weak-signal mixture detection problem that Higher
Criticism was designed to solve, with documents in place of features
and the empirical similarity null in place of an analytic Gaussian
null.

\subsection{Higher Criticism as a Cutoff Rule}
\emph{Higher Criticism} (HC) \cite{donohojin2004hc} is a
goodness-of-fit statistic designed to detect a small number of
\emph{rare and weak} signals buried in a sea of nulls --- exactly the
statistical regime of a sparse-relevance retrieval ranking. Given a
per-query score distribution and an offline null
distribution of irrelevant-passage scores, HC ranks candidates by
their $z$-score-based $p$-values and selects the cutoff that maximises
the standardised deviation between observed and expected uniform
$p$-values. The output is an adaptive $k^\star$ tied directly to the
statistical separation of relevant from irrelevant passages.

\subsection{Research Question}
\emph{Can an HC-driven cutoff match or beat fixed-$k$ and existing
adaptive baselines on answer quality while substantially reducing
context cost, and on which kinds of questions does it fail?}
This report answers that question across five datasets spanning
single-hop QA, multi-hop QA, long-context aggregation, and scientific
fact-checking, and introduces a purpose-built benchmark
(CrossEntityQA) to stress-test the sparse-relevance regime that HC
theory targets.

%====================================================================
% === Methodology ===
%====================================================================
\section{Methodology}

\subsection{HC Retrieval Algorithm}
For each query $q$ the retriever returns the top $N$ candidates with
cosine similarities $s_1 \ge s_2 \ge \dots \ge s_N$. Offline, we build
a \emph{null distribution} of similarities by sampling random
(query, non-relevant passage) pairs from the corpus. At inference each
$s_i$ is converted to a $z$-score under the null, then to a $p$-value
$p_i$. Concretely, $p_i$ is the probability of observing a similarity
\emph{at least as high as} $s_i$ between $q$ and a randomly drawn
irrelevant passage: a small $p_i$ means $s_i$ is too high to be a
noise hit, i.e.\ the candidate is unlikely to be irrelevant under
the null. Sorting $p_{(1)} \le \dots \le p_{(N)}$, the HC statistic
is
\[
\mathrm{HC}_i \;=\; \sqrt{N}\;\frac{i/N - p_{(i)}}{\sqrt{p_{(i)}(1-p_{(i)})}},
\]
which compares the empirical $i$-th order statistic $p_{(i)}$ against
its expected value $i/N$ under the null --- where the $p$-values
would be i.i.d.\ uniform on $[0,1]$, so the empirical quantiles are
the ones implied by the calibrated null distribution. A large
positive $\mathrm{HC}_i$ at some index $i$ therefore means the
top-$i$ scores are systematically more extreme than the null
predicts. The adaptive cutoff is
$k^\star = \arg\max_{i \le \gamma N} \mathrm{HC}_i$, with $\gamma$ a
small upper-fraction hyperparameter (typically $0.2$). The pipeline
returns the top-$k^\star$ passages.

\subsection{System Architecture}
Figure~\ref{fig:arch} shows the offline / online split. Offline:
chunking, embedding, FAISS indexing, and null calibration. Online:
query embedding, ANN retrieval, HC cutoff, prompt construction,
generation, and joint IR + generation evaluation.

\begin{figure}[t]
\centering
% \includegraphics[width=0.95\linewidth]{figures/fig_architecture.pdf}
\fbox{\parbox{0.9\linewidth}{\centering\vspace{6pt}
\textit{[Figure placeholder: HC-RAG architecture --- offline indexing
+ HC calibration; online retrieve, cutoff, generate, evaluate.]}\vspace{6pt}}}
\caption{HC-RAG high-level architecture.}
\label{fig:arch}
\end{figure}

\subsection{Pipelines Integrated}
We integrated HC as a drop-in cutoff stage into:
\begin{itemize}
\item \textbf{Adaptive-K} \cite{adaptivek2025} --- replaced the
largest-gap heuristic with HC at the same insertion point;
this method was designed for long-context QA in RAG.
\item \textbf{Adaptive-RAG} \cite{jeong2024adaptiverag} --- the
classifier-routed framework that dispatches each query to one of
\{no-retrieval (NOR), single-step retrieval (OneR), iterative
multi-hop retrieval (IRCoT)\}; HC was inserted to size the
retrieval at every IRCoT iteration.
\item \textbf{FlashRAG} \cite{jin2025flashrag} --- registered HC as a
new retriever post-processor compatible with their sequential,
conditional, branching, and iterative pipelines.
\item \textbf{Our pipeline} --- end-to-end RAG with configurable
HuggingFace embedders, FAISS index, multiple chunking strategies, and
OpenRouter for LLM generation; this is where all controlled HC vs.\
baseline experiments are run.
\end{itemize}

\subsection{Additional Cutoff Baselines}
For the algorithm comparison study (Section~\ref{sec:algcomp}) we
also evaluate four further selection rules that consume the same
$p$-values as HC: \textbf{Berk--Jones} \cite{berkjones1979}, which
replaces HC's numerator with a Bernoulli KL divergence;
\textbf{Bonferroni} \cite{bonferroni1961dunn} FWER control;
\textbf{Benjamini--Hochberg} \cite{benjamini1995fdr} step-up FDR
control; and \textbf{CAR} \cite{xu2025car}, a clustering-based
cut-off operating on the similarity vector directly.

\subsection{CrossEntityQA: A New Benchmark}
Most public RAG benchmarks favour single-passage answers, which masks
the regime HC targets. We constructed \textbf{CrossEntityQA}, a
synthetic QA retrieval benchmark built from Wikipedia. The construction
pipeline (Figure~\ref{fig:ceqa}) clusters Wikidata entities into
constraint-aligned groups of target size $K \in \{2, \dots, 19\}$,
attaches a Wikipedia passage per entity, and uses an LLM to write a
query that is answerable only when all $K$ passages are retrieved.
The resulting corpus is \emph{sparse}: $\sim$2\% of passages are
relevant per query, with an average of 8.8 ground-truth passages
spanning unrelated topic groups.

\begin{figure}[t]
\centering
% \includegraphics[width=0.95\linewidth]{figures/fig_ceqa_pipeline.pdf}
\fbox{\parbox{0.9\linewidth}{\centering\vspace{6pt}
\textit{[Figure placeholder: CrossEntityQA construction pipeline
(Wikidata $\to$ entity clusters $\to$ K-aligned subclusters $\to$
LLM query generation + filtering $\to$ corpus / queries / qrels).]}\vspace{6pt}}}
\caption{CrossEntityQA construction pipeline.}
\label{fig:ceqa}
\end{figure}

\subsection{Datasets}
We evaluate on five datasets spanning the regimes that stress an
adaptive cutoff: \textbf{CrossEntityQA} (our sparse multi-entity
benchmark), \textbf{HotpotQA} \cite{yang2018hotpotqa} (multi-hop QA),
\textbf{MuSiQue} \cite{trivedi2022musique} (compositional multi-hop
QA), \textbf{HoloBench} \cite{maekawa2024holobench} (long-context
database-style aggregation), and \textbf{BEIR SciFact}
\cite{thakur2021beir,wadden2020scifact} (scientific fact
verification). For failure analysis (Section~\ref{sec:failure}) we
additionally use \textbf{BEIR FiQA} \cite{thakur2021beir}.

The additional empirical studies in Section~\ref{sec:prelim}
exercise HC on three benchmarks designed in-house:
\textbf{Category-Key} (44 ``largest cities in $X$'' country queries
over 726 Wikipedia passages, with variable ground-truth size
$K \in [4,24]$); and three compound-query corpora derived from
\texttt{McAuley-Lab/Amazon-Reviews-2023} --- \textbf{Amazon
Multi-Department} (80K docs, 135 subcategory$+$price queries across
three departments), \textbf{Amazon Electronics} (100K, 44 queries),
and \textbf{Amazon Book Categories} (100K, 131 queries).

\subsection{Experimental Setup}
Embeddings use BAAI \texttt{bge-base-en-v1.5}
\cite{bgeembeddings2023} (and \texttt{bge-large-en-v1.5} for
ablations). Documents are chunked with \texttt{RecursiveChunker} at
384 tokens / 50 overlap and indexed in FAISS \cite{johnson2019faiss}.
Generation uses \texttt{gpt-4o-mini} via OpenRouter with a fixed
task-specific prompt template, except for the Amazon end-to-end
study (Section~\ref{sec:amazon}) which uses \texttt{gpt-5-mini}.
HC uses an empirical null estimated from 10K random (query,
non-relevant passage) pairs; parametric and rank-based nulls were
evaluated as alternatives, and the empirical null is retained as
default. Position-conditioned nulls remain a future-work direction
(Section~\ref{sec:future}).

\subsection{Metrics}
Retrieval: Recall@$k$, Precision@$k$, average context size (avg
$k^\star$). Generation: Exact Match (EM), token-level F1, Precision,
Recall. Efficiency: input tokens, average retrieved documents,
end-to-end retrieval latency (seconds).

%====================================================================
% === Contributions ===
%====================================================================
\section{Contributions}

\begin{enumerate}
\item \textbf{HC pipeline integrations} --- Drop-in HC cutoff in
Adaptive-K, Adaptive-RAG, FlashRAG, and a custom end-to-end pipeline.
\item \textbf{Null calibration} --- We compared empirical,
parametric, and rank-based null distributions, and retained the
empirical null as the default after observing no consistent quality
gain from the alternatives.
\item \textbf{CrossEntityQA benchmark} --- A new sparse multi-entity
retrieval dataset with variable ground-truth size ($K = 2$--$19$),
designed to stress-test adaptive-$k$ algorithms.
\item \textbf{Head-to-head evaluation} --- Five datasets, three
baselines (Top-K, Adaptive-K, Adaptive-RAG/IRCoT), reporting both
quality and cost.
\item \textbf{Negative results and root-cause analysis} --- We
characterise the two failure modes (scattered relevance, multi-hop
dependency) and a rerank-plus-threshold workaround that recovers
$+7.5\%$ recall / $+81\%$ precision on FiQA.
\end{enumerate}

%====================================================================
% === Results ===
%====================================================================
\section{Experiments and Results}

\subsection{HC vs. Top-K Baseline}
\label{sec:topk}
Table~\ref{tab:topk} summarises HC vs.\ fixed Top-K ($k=15$).
On the sparse single-hop CrossEntityQA, HC lifts EM by
$+50\%$ relative ($\CEqaEMTopK \rightarrow \CEqaEMHC$) and precision
by $+6.1\%$, at a small recall cost. On the dependency-heavy
MuSiQue, the picture inverts: HC retrieves a tighter set
(precision $\MusPrecHC$ vs.\ $\MusPrecTopK$) but loses the
cross-passage evidence chain needed to answer ($-19\%$ EM).

\begin{table}[t]
\centering
\small
\begin{tabular}{lccc}
\hline
Dataset / Metric & Top-K & HC & $\Delta$ \\
\hline
CrossEntityQA EM        & \CEqaEMTopK   & \CEqaEMHC   & $+50\%$ \\
CrossEntityQA Precision & \CEqaPrecTopK & \CEqaPrecHC & $+6.1\%$ \\
CrossEntityQA Recall    & \CEqaRecTopK  & \CEqaRecHC  & $-6.5\%$ \\
\hline
MuSiQue EM              & \MusEMTopK    & \MusEMHC    & $-19.2\%$ \\
MuSiQue Precision       & \MusPrecTopK  & \MusPrecHC  & $+12.7\%$ \\
MuSiQue Recall          & \MusRecTopK   & \MusRecHC   & $-26.4\%$ \\
MuSiQue F1              & \MusFOneTopK  & \MusFOneHC  & $-10.8\%$ \\
\hline
\end{tabular}
\caption{HC vs.\ Top-K ($k=15$). HC wins on sparse single-hop QA,
loses on dependency-heavy multi-hop QA.}
\label{tab:topk}
\end{table}

\subsection{HC vs. Adaptive-K}
\label{sec:adak}
On HoloBench (long-context database-style aggregation) HC dominates
Adaptive-K on every quality metric while using fewer documents and
tokens: precision $\HolPrecHC$ vs.\ $\HolPrecAdaK$ ($+67.6\%$
relative), F1 $\HolFOneHC$ vs.\ $\HolFOneAdaK$ ($+42.1\%$), EM
$\HolEMHC$ vs.\ $\HolEMAdaK$ ($+14.8\%$), with a small recall
trade-off ($\HolRecHC$ vs.\ $\HolRecAdaK$). On HotpotQA the precision
gap is even larger: HC doubles Adaptive-K's precision ($\HotPrecHC$
vs.\ $\HotPrecAdaK$; $+100\%$ relative) and lifts F1 by $+60\%$
($\HotFOneHC$ vs.\ $\HotFOneAdaK$) and EM by $+23\%$ ($\HotEMHC$
vs.\ $\HotEMAdaK$). On BEIR SciFact HC adds $+10\%$ precision and
$+4\%$ recall over Adaptive-K.

\subsection{HC vs. Adaptive-RAG (IRCoT)}
Against the multi-step Adaptive-RAG router on HotpotQA, HC
underperforms by a uniform single-digit margin: EM $\HotEMHC$ vs.\
$\HotEMIRCoT$ ($-6.7\%$), precision $\HotPrecHC$ vs.\
$\HotPrecIRCoT$ ($-9.9\%$), F1 $\HotFOneHC$ vs.\ $\HotFOneIRCoT$
($-10.4\%$). IRCoT's iterative query reformulation recovers
cross-document evidence that a single-shot HC cutoff cannot, at the
cost discussed below.

\subsection{Efficiency}
Table~\ref{tab:eff} shows the cost side. On HotpotQA HC uses
$94.77\%$ fewer input tokens ($\HotTokHC$ vs.\ $\HotTokAdaK$),
retrieves $\sim$21$\times$ fewer documents ($\HotDocsHC$ vs.\
$\HotDocsAdaK$), and is $90.25\%$ faster ($\HotLatHC$\,s vs.\
$\HotLatAdaK$\,s end-to-end). On HoloBench the gains carry over:
$71\%$ fewer tokens ($\HolTokHC$ vs.\ $\HolTokAdaK$), $\sim$3.6$\times$
fewer documents ($\HolDocsHC$ vs.\ $\HolDocsAdaK$), and $\sim$8$\times$
faster ($\HolLatHC$\,s vs.\ $\HolLatAdaK$\,s). This is the central
practical pitch for HC: when it matches or beats baselines on quality,
it does so at a fraction of the dollar and latency budget.

\begin{table}[t]
\centering
\small
\begin{tabular}{lccc}
\hline
Dataset & Metric & HC & Adaptive-K \\
\hline
HotpotQA  & Input tokens     & \HotTokHC  & \HotTokAdaK \\
HotpotQA  & Avg \# docs      & \HotDocsHC & \HotDocsAdaK \\
HotpotQA  & Retrieval lat.\ (s) & \HotLatHC  & \HotLatAdaK \\
\hline
HoloBench & Input tokens     & \HolTokHC  & \HolTokAdaK \\
HoloBench & Avg \# docs      & \HolDocsHC & \HolDocsAdaK \\
HoloBench & Retrieval lat.\ (s) & \HolLatHC  & \HolLatAdaK \\
\hline
\end{tabular}
\caption{Efficiency: HC vs.\ Adaptive-K.}
\label{tab:eff}
\end{table}

\subsection{Failure Analysis}
\label{sec:failure}
To understand when HC fails we ran a focused position study on BEIR
FiQA. HC implicitly assumes that relevant passages \emph{cluster} at
the top of the ranked list; the test then detects the boundary where
the score distribution stops being unusual. On FiQA the relevant
passages are instead \emph{scattered}: only $6/100$ queries have
$\ge 2$ consecutive relevant passages at the top, and the median
position of a relevant passage is 5. Even after cross-encoder
reranking (which lifts the bi-encoder Cohen's $d$ from $0.01$ to
$1.13$) clustering improves only marginally (to $8/100$).

A simple rerank-plus-fixed-threshold rule recovers most of what HC
loses: \texttt{score $\ge 0$} after cross-encoder rerank delivers
$+7.5\%$ recall and $+81\%$ precision over the Top-5 baseline,
suggesting a productive hybrid (Section~\ref{sec:future}).

%====================================================================
% === Additional Empirical Studies ===
%====================================================================
\section{Additional Empirical Studies}
\label{sec:prelim}

This section reports three studies that broaden the empirical
picture beyond the long-form QA benchmarks above: a controlled-$K$
benchmark designed to isolate HC's adaptivity, a large-corpus
compound-query retrieval study with end-to-end generation, and a
head-to-head comparison of HC against five alternative adaptive
cut-off rules.

\subsection{Controlled-$K$ Benchmark: Category-Key}
On Category-Key (44 ``largest cities in $X$'' queries over 726
Wikipedia passages; BGE bi-encoder, FAISS), HC achieves F1
$\CatHCF$ (Recall $\CatHCR$, Precision $\CatHCP$) at mean
$k^\star = \CatHCK$, beating every fixed top-$k$ in the sweep
$k \in \{5,8,10,12,15,20\}$ (best fixed F1 $= 0.783$ at $k{=}10$).
The Pearson correlation between HC's adaptive $k^\star$ and the
per-query ground-truth $K$ is $\CatCorrK$ --- direct evidence that
the HC argmax tracks the true relevance boundary rather than
collapsing to a constant.

\subsection{Compound-Query Retrieval: Amazon (80K--100K)}
\label{sec:amazon}
The Amazon benchmarks stress HC on heterogeneous product
embeddings (OpenAI \texttt{text-embedding-3-small}, FAISS) and
compound subcategory$+$price queries that standard embeddings
encode only weakly. On the 80K \textbf{Multi-Department} corpus
(135 queries spanning Clothing, Electronics, Home \& Kitchen), HC
reaches retrieval F1 $\AmMdHCF$ at mean $k^\star = \AmMdHCK$,
ahead of Top-20 ($\AmMdTwentyF$; $+9.7\%$), Top-29 set to the
average number of relevant documents per query ($\AmMdAvgF$;
$+6.3\%$), and Top-50 ($\AmMdFiftyF$; $+8.9\%$), while retrieving
$32\%$ fewer documents than Top-50. End-to-end with a
\texttt{gpt-5-mini} generator HC posts Token-F1 $\AmMdHCTokF$
vs.\ $\AmMdTwentyTokF$ for Top-20 ($+14.8\%$), with an LLM-judge
Answer Faithfulness near $5.0$ (the maximum) across all methods,
indicating the generator already filters distractors effectively.
The same pattern holds on \textbf{Electronics} (HC F1 $\AmElHCF$
vs.\ Top-20's $\AmElTwentyF$; $+3.5\%$) and \textbf{Book
Categories} (HC F1 $\AmBkHCF$ vs.\ Top-50's $\AmBkFiftyF$;
$+4.0\%$, with $37\%$ fewer documents).

\subsection{Algorithm Comparison: HC vs.\ Other Adaptive Cut-offs}
\label{sec:algcomp}
On the same 135 Amazon Multi-Department queries we compared HC
head-to-head against five further adaptive cut-off rules
(Table~\ref{tab:algcomp}), all sharing an identical FAISS pool of
$N{=}1000$ candidates. HC wins on F1 ($\AlgHCF$), with
Berk--Jones a close second ($\AlgBJF$); the gap-based Adaptive-K
\cite{adaptivek2025} sits well behind ($\AlgAdaKF$); the
clustering-based CAR \cite{xu2025car} collapses at its default
pool ($\AlgCARF$) and matches HC only after pool retuning. The
classical multiple-testing rules under-retrieve drastically:
Bonferroni ($\AlgBonfF$) and Benjamini--Hochberg ($\AlgBHF$)
reject almost no candidates at a flat significance threshold,
confirming that an \emph{adaptive} test statistic (HC, BJ) is
essential when per-query signal strength varies.

\begin{table}[t]
\centering
\small
\begin{tabular}{lcc}
\hline
Method & F1 & Mean $k$ \\
\hline
HC + rescale     & $\AlgHCF$    & 34.0 \\
Berk--Jones      & $\AlgBJF$    & 34.8 \\
Top-50           & 0.187        & 50.0 \\
Top-20           & 0.185        & 20.0 \\
Adaptive-K (gap) & $\AlgAdaKF$  & 9.0 \\
CAR (cluster)    & $\AlgCARF$   & 416.5 \\
BH               & $\AlgBHF$    & 1.6 \\
Bonferroni       & $\AlgBonfF$  & 0.3 \\
\hline
\end{tabular}
\caption{Adaptive cut-off comparison on Amazon Multi-Department
(135 compound queries, pool $N{=}1000$). HC posts the best F1;
flat-threshold rules under-retrieve.}
\label{tab:algcomp}
\end{table}

%====================================================================
% === Discussion ===
%====================================================================
\section{Discussion}

\paragraph{When HC wins.} Sparse-relevance corpora with
well-separated embeddings and clustered relevant passages
(CrossEntityQA, HoloBench, BEIR SciFact, Category-Key, Amazon
Multi-Department). Here HC's adaptivity translates directly into
precision gains and large token savings.

\paragraph{When HC loses.} Two regimes. (i) \emph{Dependency-heavy
multi-hop} questions (MuSiQue, HotpotQA vs.\ IRCoT) where the answer
requires evidence that an iterative reformulation can find but a
single cutoff cannot. (ii) \emph{Scattered relevance} corpora
(FiQA) where the clustering assumption is violated.

\paragraph{Limitations.} Across the regimes where HC wins it tends to
lift Precision and lower Recall versus fixed Top-K. We read this as
the expected finite-$N$ behaviour of
$k^\star = \arg\max_{i \le \gamma N} \mathrm{HC}_i$ in the
rare-and-weak regime: the standardised deviation peaks at the
smallest $i$ for which the empirical $p_{(i)}$ already sits well
below the uniform expectation $i/N$, so the rule prefers tight
cutoffs whose entries are extreme under the null --- a precision-leaning
operating point with a recall trade-off, rather than a contradiction
of the underlying theory. We did not turn this into a formal
finite-$N$ claim, and we ran the per-(query, non-relevant) $p$-value
uniformity diagnostic only on FiQA, where we also prototyped
position-aware and gap-based alternative nulls; running the same
diagnostic across the full benchmark suite is the natural empirical
companion (Section~\ref{sec:future}).

\paragraph{Response to presentation feedback.}
Feedback from the final presentation covered conceptual,
methodological, and empirical points, which are addressed throughout
this report: (i)~retrieval is framed up front as a sparse
signal-detection problem (Section~\ref{sec:sparse}) so that the
reason HC applies at all is explicit before the algorithm is defined;
(ii)~the HC Retrieval Algorithm subsection now spells out, in plain
language, what the $(q,d)$ $p$-value is, that the HC statistic
compares empirical order-statistic quantiles against the uniform null
implied by the calibrated null distribution, and how this drives the
adaptive cutoff; (iii)~null-calibration choices are made explicit in
the Methodology, with the empirical null retained as the default
after comparison against parametric and rank-based alternatives;
position-aware and gap-based variants were prototyped on FiQA, and
a unified $p$-value uniformity diagnostic across all evaluation
datasets is left for future work;
(iv)~the failure-mode taxonomy (scattered relevance vs.\ multi-hop
dependency) is stated directly in the abstract and revisited in this
section; (v)~the rerank-plus-threshold hybrid is reported as an
explicit comparison point in Section~\ref{sec:failure} rather than
only as future work; and (vi)~efficiency numbers (tokens, document
counts, latency) are reported alongside every quality result so that
gains and trade-offs are directly comparable.

%====================================================================
% === Conclusion ===
%====================================================================
\section{Conclusions and Future Work}
\label{sec:future}

We showed that Higher Criticism is a competitive query-adaptive
cutoff for RAG: matching or beating fixed Top-K and Adaptive-K on
four of five long-form QA datasets while reducing input-token cost
by up to $\sim$95\%; on the controlled Category-Key benchmark HC
also beats every fixed top-$k$ in the sweep, and on the Amazon
compound-query benchmark it leads five alternative adaptive
cut-off rules including Berk--Jones, BH, Bonferroni, gap-based
Adaptive-K, and clustering-based CAR. We documented the two
regimes where HC fails. Three directions stand out for future
work:
\begin{itemize}
\item \textbf{Rerank $+$ HC hybrid.} Apply HC after a cross-encoder
rerank to address the scattered-relevance failure mode; the FiQA
threshold experiment is a positive existence proof.
\item \textbf{Iterative HC.} Embed HC inside Adaptive-RAG's IRCoT
loop to combine multi-hop coverage with adaptive-cost control.
\item \textbf{Position-aware nulls.} Replace the per-query
empirical null with a position-conditioned null to reduce variance
without destroying HC's signal-detection power.
\item \textbf{Per-dataset $p$-value uniformity diagnostic.} Run the
FiQA-style KS uniformity check on the calibrated $p$-values across
all evaluation datasets, paired with a finite-$N$ analysis of
$\arg\max \mathrm{HC}_i$, to convert the precision-leaning operating
point from an empirical observation into a falsifiable prediction.
\end{itemize}

%====================================================================
% === References ===
%====================================================================
\bibliographystyle{aaai}
\bibliography{refs}

%====================================================================
% === Appendix Roadmap ===
% Appendices do not count toward the 6-page body limit.
%====================================================================
\clearpage
\onecolumn
\section*{Appendix A: Project Roadmap}

\begin{center}
\small
\begin{tabular}{p{0.07\linewidth}p{0.88\linewidth}}
\hline
Week & Milestone / Outcome \\
\hline
4   & \textbf{Foundations.} Implemented the HC module
($p$-values from a per-query null, $\gamma$ upper-fraction cutoff)
and built the indexing stack on the CRAG \cite{crag2024} corpus
(\texttt{RecursiveChunker}, 384/50 overlap; LLM interface over
OpenRouter). \emph{Direction change}: a small $k$-sweep with
\texttt{all-mpnet-base-v2} showed $k{=}10$ as the best fixed
operating point, but switching to BAAI \texttt{bge-base-en-v1.5}
\cite{bgeembeddings2023} on a 1K-doc subset gave a meaningful
recall lift, so BGE was adopted as the default embedder for the
rest of the project. \\
5   & \textbf{Survey of alternative adaptive retrievers.} Read and
catalogued three competing approaches --- Reciprocal Rank Fusion,
Cluster-Adaptive Retrieval (CAR) \cite{xu2025car}, and
uncertainty-based filtering --- to position HC as a global,
distribution-aware test rather than a per-document heuristic, and
locked in the comparison set for later head-to-heads (Bonferroni
\cite{bonferroni1961dunn}, Benjamini--Hochberg
\cite{benjamini1995fdr}, gap-based Adaptive-K
\cite{adaptivek2025}). First HC vs.\ Top-K runs on our pipeline. \\
6   & \textbf{Long-context regime: HoloBench.} First head-to-head
HC vs.\ gap-based Adaptive-K on HoloBench
\cite{maekawa2024holobench}. Quality scores tracked closely on the
3-level LLM judge --- \emph{iteration}: diagnosed as a
generation/judge bottleneck on database-style aggregation, not a
retrieval-side tie, which motivated stronger retrieval-only metrics
in later rounds. \\
7   & \textbf{Multi-hop regime + dataset-aware $\gamma$.} Repeated
the comparison on HotpotQA \cite{yang2018hotpotqa}. \emph{Change in
direction}: data analysis exposed a stark dense vs.\ sparse
asymmetry between HotpotQA and HoloBench, and the \emph{same}
$\gamma$ behaved very differently on the two ---
$\gamma$ was reframed as a dataset-aware hyperparameter rather than
a single global value. Began integration work for Adaptive-RAG
\cite{jeong2024adaptiverag} and FlashRAG \cite{jin2025flashrag}. \\
8   & \textbf{Pipeline integrations + FiQA failure.} Two parallel
tracks. (a)~Plugged HC into Adaptive-RAG as a fourth strategy
alongside NOR / OneR / IRCoT; on HotpotQA HC beat single-step OneR
and lost to multi-step IRCoT, as expected on a multi-hop benchmark.
(b)~Registered HC and Berk--Jones \cite{berkjones1979} as
FlashRAG / LlamaIndex post-processors and ran them against
gap-based Adaptive-K and Fixed Top-$10$ on dev subsets of BEIR
SciFact \cite{thakur2021beir,wadden2020scifact}, BEIR FiQA
\cite{thakur2021beir}, and MS MARCO. \emph{Iteration}: FiQA was
the first clear single-hop negative result and triggered a position
study, which in turn produced a rerank-plus-fixed-threshold
workaround that recovered most of the lost recall and precision
(Section~\ref{sec:failure}). \\
9   & \textbf{CrossEntityQA construction.} \emph{Change in
direction}: rather than continue tuning HC against off-the-shelf
benchmarks, we built a purpose-designed sparse-relevance benchmark
to surface the regime HC theory targets --- SPARQL-driven Wikidata
entity clusters $\to$ Wikipedia passages $\to$ LLM-generated queries
that are answerable only by aggregating multiple passages, with
$\sim$2\% relevance and a controllable $K$. \\
10  & \textbf{Negative result on CrossEntityQA + BM25.} First
CrossEntityQA evaluation went through Adaptive-RAG with BM25 as the
retriever and HC \emph{lost} to fixed Top-$5$. Root cause: BM25
score magnitude encodes lexical match, not the temporal /
categorical / geographic constraints these queries impose, so HC's
test was fed a misaligned signal --- it over-selected on
entity-saturated queries and under-selected on generic constraint
queries. \emph{Change in direction}: dropped BM25 for
CrossEntityQA and committed to a dense retriever for the final
runs. \\
11  & \textbf{Final head-to-head.} Re-ran HC vs.\ Top-K on
CrossEntityQA with the dense pipeline --- HC lifted EM by $\sim$50\%
relative (Table~\ref{tab:topk}) --- and replicated the multi-hop
failure mode on MuSiQue \cite{trivedi2022musique}, closing the
loop on the ``when HC wins / when HC loses'' picture used in this
report. \\
\hline
\end{tabular}
\end{center}

\end{document}
